\section{Experimental Environment}
\label{sec:methodology_experimental_environment}

\subsection{Benchmark Suite}
\label{subsec:methodology_benchmark_suite}

The benchmark suite consisted of five standard reinforcement-learning control environments: \textsc{MountainCar-v0}, \textsc{BipedalWalker-v3}, \textsc{Ant-v5}, \textsc{HalfCheetah-v5}, and \textsc{Humanoid-v5}.
The suite was selected to cover both sparse-reward and continuous-control settings.
\textsc{MountainCar-v0} provides a sparse and delayed success signal in a low-dimensional discrete-action task.
\textsc{BipedalWalker-v3} requires coordinated continuous control and contains failure states.
The three MuJoCo locomotion tasks provide higher-dimensional continuous-control settings with denser reward feedback and different morphology complexities.

The default reward functions and termination conditions were used in all environments.
No domain randomization was applied.
No additional frame skipping was introduced beyond the environment configuration.
Training used synchronous vectorized environments, with the number of parallel instances selected to obtain a nominal PPO batch size of $N=16{,}384$ transitions per update.

\begin{table}[!htbp]
\centering
\small
\begin{fitwidthtabular}{lrrrrr}
\hline
Environment & $B$ & $T$ & $N$ & Total steps & Seeds \\
\hline
\textsc{MountainCar-v0}   & 16 & 1{,}024 & 16{,}384 & 3{,}000{,}000  & 10 \\
\textsc{BipedalWalker-v3} &  8 & 2{,}048 & 16{,}384 & 7{,}000{,}000  &  8 \\
\textsc{Ant-v5}           &  8 & 2{,}048 & 16{,}384 & 15{,}000{,}000 &  8 \\
\textsc{HalfCheetah-v5}   &  8 & 2{,}048 & 16{,}384 & 15{,}000{,}000 &  8 \\
\textsc{Humanoid-v5}      &  4 & 4{,}096 & 16{,}384 & 30{,}000{,}000 &  5 \\
\hline
\end{fitwidthtabular}
\caption[Benchmark suite and training budgets]{Benchmark suite and training budgets. $B$ is the number of parallel environment instances, $T$ is the rollout horizon per instance, and $N=B\times T$ is the nominal PPO update batch size.}
\label{tab:methodology_env_suite}
\end{table}

\subsection{Training Budgets}
\label{subsec:methodology_training_budgets}

The total environment-step budgets in Table~\ref{tab:methodology_env_suite} were chosen to extend beyond the early learning phase.
This allowed the evaluation to measure both rapid initial improvement and long-horizon stability.
The final PPO update in a run was allowed to contain fewer than $N$ transitions if the remaining step budget was smaller than a full rollout batch.

The GLPE hyperparameters used for each environment are shown in Table~\ref{tab:methodology_glpe_hparams}.
The same intrinsic scaling coefficient $\eta$ and intrinsic clipping range $r_{\max}$ were used by all intrinsic methods within a given environment.
This kept the reward-scale comparison consistent across ICM, RND, RIDE, RIAC, GLPE, and GLPE (no gate).

\begin{table}[!htbp]
\centering
\small
\begin{fitwidthtabular}{lrrrrrrrrrrrrr}
\hline
Environment & $\eta$ & $r_{\max}$ & $\alpha_{\mathrm{LP}}$ & $p_{\mathrm{start}}$ & $p_{\mathrm{end}}$ & $C$ & $D_{\max}$ & $\beta_{\mathrm{long}}$ & $\beta_{\mathrm{short}}$ & $\kappa$ & $\tau_s$ & $K$ & $R_{\min}$ \\
\hline
\textsc{MountainCar-v0}   & 0.05 & 4.0 & 0.5 & 0.05 & 0.80 & 128 & 10 & 0.995 & 0.90 & 0.010 & 2.0 & 5 &  8 \\
\textsc{BipedalWalker-v3} & 0.08 & 4.0 & 0.5 & 0.12 & 0.75 & 200 & 12 & 0.995 & 0.90 & 0.010 & 2.0 & 5 & 16 \\
\textsc{Ant-v5}           & 0.06 & 4.0 & 0.6 & 0.10 & 0.70 & 256 & 12 & 0.997 & 0.92 & 0.012 & 2.5 & 8 & 64 \\
\textsc{HalfCheetah-v5}   & 0.04 & 5.0 & 0.5 & 0.05 & 0.75 & 256 & 12 & 0.995 & 0.90 & 0.010 & 2.0 & 5 & 32 \\
\textsc{Humanoid-v5}      & 0.06 & 5.0 & 0.5 & 0.05 & 0.80 & 256 & 12 & 0.998 & 0.92 & 0.010 & 2.0 & 5 & 32 \\
\hline
\end{fitwidthtabular}
\caption[GLPE hyperparameters]{GLPE hyperparameters by environment. The impact weight was fixed at $\alpha_{\mathrm{impact}}=1.0$, and the hysteresis multiplier was fixed at $h=2.0$. GLPE (no gate) used the same settings but disabled the gate.}
\label{tab:methodology_glpe_hparams}
\end{table}

\subsection{Random Seeds and Determinism}
\label{subsec:methodology_random_seeds_and_determinism}

Each method was evaluated using the same seed set within each environment.
The seed counts are shown in Table~\ref{tab:methodology_env_suite}.
The experiments enabled deterministic execution where supported by the software stack and environment backend.
The purpose of fixed seed sets was not to remove randomness from reinforcement learning, but to ensure that method comparisons were paired as closely as practical across independent training runs.